% --- Archon physics formalization source begin ---
% archon:physics
% archon:covers PhyXMiniProblems/problem_phyx_mini_0554.lean
% archon:source-report reports/phyx_mini/problem_phyx_mini_0554.source.json

\chapter{Physics problem phyx\_mini\_0554}
\label{ch:PhyXMiniProblems_problem_phyx_mini_0554}

\paragraph{Problem source.}
\#\# Physical scenario

In the reaction \textbackslash{}( \textbackslash{}text\{K\}\textasciicircum{}- + \textbackslash{}text\{p\} \textbackslash{}rightarrow \textbackslash{}Lambda\textasciicircum{}0 + \textbackslash{}pi\textasciicircum{}0 \textbackslash{}), a charged K meson (mass \textbackslash{}( 493.7 \textbackslash{}text\{ MeV\}/c\textasciicircum{}2 \textbackslash{})) collides with a proton (\textbackslash{}( 938.3 \textbackslash{}text\{ MeV\}/c\textasciicircum{}2 \textbackslash{})) at rest, producing a lambda particle (\textbackslash{}( 1115.7 \textbackslash{}text\{ MeV\}/c\textasciicircum{}2 \textbackslash{})) and a neutral pi meson (\textbackslash{}( 135.0 \textbackslash{}text\{ MeV\}/c\textasciicircum{}2 \textbackslash{})), as represented in figure. The initial kinetic energy of the K meson is \textbackslash{}( 152.4 \textbackslash{}text\{ MeV\} \textbackslash{}). After the interaction, the pi meson has a kinetic energy of \textbackslash{}( 254.8 \textbackslash{}text\{ MeV\} \textbackslash{}).

\#\# Question

Find the kinetic energy of the lambda.

\#\# Answer choices

- A: A: \textbackslash{}( 65.6 \textbackslash{}text\{ MeV\} \textbackslash{})

- B: B: \textbackslash{}( 90.4 \textbackslash{}text\{ MeV\} \textbackslash{})

- C: C: \textbackslash{}( 33.1 \textbackslash{}text\{ MeV\} \textbackslash{})

- D: D: \textbackslash{}( 78.9 \textbackslash{}text\{ MeV\} \textbackslash{})

\#\# Figure caption (auxiliary; use the image as primary evidence)

The image consists of two panels labeled (a) and (b), each depicting a physics diagram with particles and coordinate axes.

**Panel (a):**
- It shows two particles on a 2D coordinate system with axes labeled \textbackslash{}(x\textbackslash{}) and \textbackslash{}(y\textbackslash{}).
- A particle labeled \textbackslash{}(K\textasciicircum{}-\textbackslash{}) is moving along the negative x-axis towards a stationary particle labeled \textbackslash{}(p\textbackslash{}).
- There is an arrow indicating the direction of motion.

**Panel (b):**
- It depicts another scene on a 2D coordinate system with axes \textbackslash{}(x\textbackslash{}) and \textbackslash{}(y\textbackslash{}).
- Two particles are shown: \textbackslash{}(\textbackslash{}pi\textasciicircum{}0\textbackslash{}) and \textbackslash{}(\textbackslash{}Lambda\textasciicircum{}0\textbackslash{}).
- The particle \textbackslash{}(\textbackslash{}pi\textasciicircum{}0\textbackslash{}) is moving at an angle \textbackslash{}(\textbackslash{}theta\textbackslash{}) from the y-axis, depicted by an arrow.
- The particle \textbackslash{}(\textbackslash{}Lambda\textasciicircum{}0\textbackslash{}) is depicted in motion at an angle \textbackslash{}(\textbackslash{}phi\textbackslash{}) from the x-axis, shown by another arrow.
- The angles \textbackslash{}(\textbackslash{}theta\textbackslash{}) and \textbackslash{}(\textbackslash{}phi\textbackslash{}) are marked with curved arrows.

Both panels seem to illustrate some interaction or decay process involving particles with labeled directions.

\#\# Dataset metadata

Particle Physics; Multi-Formula Reasoning, Physical Model Grounding Reasoning

\paragraph{Recorded answer/context.}
D: D: \textbackslash{}( 78.9 \textbackslash{}text\{ MeV\} \textbackslash{})

\paragraph{Figure/image path.}
/root/proposal\_for\_physic/science-mango/phyx\_mini\_run/phyx\_data/test\_image/554.png

\paragraph{Formalization target.}
create a compiling Lean file with sorry bodies at `PhyXMiniProblems/problem\_phyx\_mini\_0554.lean`. The Lean declarations must preserve the physical quantities, dimensions or dimensional roles, figure labels, governing-law hypotheses, and final relation expressed by this problem.
Use Mathlib/Physlib names found through LeanExplore where available. If a physics API is missing, introduce faithful local abstractions rather than scalar placeholder aliases.

\begin{theorem}[Physics formalization target]
\label{thm:physics:phyx_mini_0554:target}
\lean{PhyXMiniProblems.ProblemPhyXMini0554.problem_phyx_mini_0554}
\uses{def:physics:phyx-mini-0554:phyxminiproblems-problemphyxmini0554-energyinmegaelectronvolts, def:physics:phyx-mini-0554:phyxminiproblems-problemphyxmini0554-kaonprotonreactionsetup, def:physics:phyx-mini-0554:phyxminiproblems-problemphyxmini0554-matchesreactionproblemdata, def:physics:phyx-mini-0554:phyxminiproblems-problemphyxmini0554-matchesprimaryreactionfigure, def:physics:phyx-mini-0554:phyxminiproblems-problemphyxmini0554-hasphysicalreactionparameters, def:physics:phyx-mini-0554:phyxminiproblems-problemphyxmini0554-satisfiesrelativisticreactionlaws, def:physics:phyx-mini-0554:phyxminiproblems-problemphyxmini0554-displayedlambdakineticenergymev, def:physics:phyx-mini-0554:phyxminiproblems-problemphyxmini0554-recordedanswerchoice}
The assigned autoformalize agent should translate this physics problem into Lean declarations in the covered file, with theorem and lemma proof bodies written as `by sorry`.
\end{theorem}
\begin{proof}
This is an autoformalization task, not a proof task. Produce statements that can later be proved by the physics prover without weakening the physical model.
\end{proof}

% --- Archon declaration topology begin ---
\section{Declaration topology}
\begin{definition}[Rest Mass Quantity]
  \label{def:physics:phyx-mini-0554:phyxminiproblems-problemphyxmini0554-restmassquantity}
  \lean{PhyXMiniProblems.ProblemPhyXMini0554.RestMassQuantity}
  A nonnegative physical rest mass with mass dimension M
\end{definition}
\begin{proof}
  This is the declared model definition of Rest Mass Quantity.
\end{proof}
\begin{definition}[Plane Momentum]
  \label{def:physics:phyx-mini-0554:phyxminiproblems-problemphyxmini0554-planemomentum}
  \lean{PhyXMiniProblems.ProblemPhyXMini0554.PlaneMomentum}
  A dimensionful two-dimensional momentum in the plane of the diagram
\end{definition}
\begin{proof}
  This is the declared model definition of Plane Momentum.
\end{proof}
\begin{definition}[Diagram Axis]
  \label{def:physics:phyx-mini-0554:phyxminiproblems-problemphyxmini0554-diagramaxis}
  \lean{PhyXMiniProblems.ProblemPhyXMini0554.DiagramAxis}
  The two coordinate axes printed in both panels of the primary image
\end{definition}
\begin{proof}
  This is the declared model definition of Diagram Axis.
\end{proof}
\begin{definition}[to Fin]
  \label{def:physics:phyx-mini-0554:phyxminiproblems-problemphyxmini0554-diagramaxis-tofin}
  \lean{PhyXMiniProblems.ProblemPhyXMini0554.DiagramAxis.toFin}
  \uses{def:physics:phyx-mini-0554:phyxminiproblems-problemphyxmini0554-diagramaxis}
  The Fin 2 coordinate corresponding to a diagram-axis label
\end{definition}
\begin{proof}
  This is the declared model definition of to Fin.
\end{proof}
\begin{definition}[rest Mass In Kilograms]
  \label{def:physics:phyx-mini-0554:phyxminiproblems-problemphyxmini0554-restmassinkilograms}
  \lean{PhyXMiniProblems.ProblemPhyXMini0554.restMassInKilograms}
  \uses{def:physics:phyx-mini-0554:phyxminiproblems-problemphyxmini0554-restmassquantity}
  Kilogram readout of a physical rest mass
\end{definition}
\begin{proof}
  This is the declared model definition of rest Mass In Kilograms.
\end{proof}
\begin{definition}[energy In Joules]
  \label{def:physics:phyx-mini-0554:phyxminiproblems-problemphyxmini0554-energyinjoules}
  \lean{PhyXMiniProblems.ProblemPhyXMini0554.energyInJoules}
  Joule readout of a physical energy
\end{definition}
\begin{proof}
  This is the declared model definition of energy In Joules.
\end{proof}
\begin{definition}[energy In Mega Electron Volts]
  \label{def:physics:phyx-mini-0554:phyxminiproblems-problemphyxmini0554-energyinmegaelectronvolts}
  \lean{PhyXMiniProblems.ProblemPhyXMini0554.energyInMegaElectronVolts}
  \uses{def:physics:phyx-mini-0554:phyxminiproblems-problemphyxmini0554-energyinjoules}
  MeV readout of a physical energy
\end{definition}
\begin{proof}
  This is the declared model definition of energy In Mega Electron Volts.
\end{proof}
\begin{definition}[speed Of Light In Meters Per Second]
  \label{def:physics:phyx-mini-0554:phyxminiproblems-problemphyxmini0554-speedoflightinmeterspersecond}
  \lean{PhyXMiniProblems.ProblemPhyXMini0554.speedOfLightInMetersPerSecond}
  Physlib's exact vacuum speed of light, read in metres per second
\end{definition}
\begin{proof}
  This is the declared model definition of speed Of Light In Meters Per Second.
\end{proof}
\begin{definition}[rest Energy Equivalent In Joules]
  \label{def:physics:phyx-mini-0554:phyxminiproblems-problemphyxmini0554-restenergyequivalentinjoules}
  \lean{PhyXMiniProblems.ProblemPhyXMini0554.restEnergyEquivalentInJoules}
  \uses{def:physics:phyx-mini-0554:phyxminiproblems-problemphyxmini0554-restmassquantity, def:physics:phyx-mini-0554:phyxminiproblems-problemphyxmini0554-restmassinkilograms, def:physics:phyx-mini-0554:phyxminiproblems-problemphyxmini0554-speedoflightinmeterspersecond}
  Rest-energy equivalent m c\textasciicircum{}2, expressed in joules
\end{definition}
\begin{proof}
  This is the declared model definition of rest Energy Equivalent In Joules.
\end{proof}
\begin{definition}[rest Mass In Mega Electron Volts Per C2]
  \label{def:physics:phyx-mini-0554:phyxminiproblems-problemphyxmini0554-restmassinmegaelectronvoltsperc2}
  \lean{PhyXMiniProblems.ProblemPhyXMini0554.restMassInMegaElectronVoltsPerC2}
  \uses{def:physics:phyx-mini-0554:phyxminiproblems-problemphyxmini0554-restmassquantity, def:physics:phyx-mini-0554:phyxminiproblems-problemphyxmini0554-energyinjoules, def:physics:phyx-mini-0554:phyxminiproblems-problemphyxmini0554-restenergyequivalentinjoules}
  Read a rest mass in the high-energy unit MeV/c\textasciicircum{}2. Numerically this is the mass's rest-energy equivalent divided by one MeV
\end{definition}
\begin{proof}
  This is the declared model definition of rest Mass In Mega Electron Volts Per C2.
\end{proof}
\begin{definition}[momentum Component In SI]
  \label{def:physics:phyx-mini-0554:phyxminiproblems-problemphyxmini0554-momentumcomponentinsi}
  \lean{PhyXMiniProblems.ProblemPhyXMini0554.momentumComponentInSI}
  \uses{def:physics:phyx-mini-0554:phyxminiproblems-problemphyxmini0554-planemomentum, def:physics:phyx-mini-0554:phyxminiproblems-problemphyxmini0554-diagramaxis}
  One momentum component read in coherent SI units kg m / s
\end{definition}
\begin{proof}
  This is the declared model definition of momentum Component In SI.
\end{proof}
\begin{definition}[momentum Magnitude In SI]
  \label{def:physics:phyx-mini-0554:phyxminiproblems-problemphyxmini0554-momentummagnitudeinsi}
  \lean{PhyXMiniProblems.ProblemPhyXMini0554.momentumMagnitudeInSI}
  \uses{def:physics:phyx-mini-0554:phyxminiproblems-problemphyxmini0554-planemomentum, def:physics:phyx-mini-0554:phyxminiproblems-problemphyxmini0554-momentumcomponentinsi}
  Euclidean magnitude of a plane momentum in kg m / s
\end{definition}
\begin{proof}
  This is the declared model definition of momentum Magnitude In SI.
\end{proof}
\begin{definition}[Particle Label]
  \label{def:physics:phyx-mini-0554:phyxminiproblems-problemphyxmini0554-particlelabel}
  \lean{PhyXMiniProblems.ProblemPhyXMini0554.ParticleLabel}
  The four particle species named in the reaction channel
\end{definition}
\begin{proof}
  This is the declared model definition of Particle Label.
\end{proof}
\begin{definition}[Reaction Side]
  \label{def:physics:phyx-mini-0554:phyxminiproblems-problemphyxmini0554-reactionside}
  \lean{PhyXMiniProblems.ProblemPhyXMini0554.ReactionSide}
  Whether a particle belongs to the incoming or outgoing side
\end{definition}
\begin{proof}
  This is the declared model definition of Reaction Side.
\end{proof}
\begin{definition}[particle Side]
  \label{def:physics:phyx-mini-0554:phyxminiproblems-problemphyxmini0554-particleside}
  \lean{PhyXMiniProblems.ProblemPhyXMini0554.particleSide}
  \uses{def:physics:phyx-mini-0554:phyxminiproblems-problemphyxmini0554-particlelabel, def:physics:phyx-mini-0554:phyxminiproblems-problemphyxmini0554-reactionside}
  The reaction side fixed by K\textasciicircum{}- + p -> Lambda\textasciicircum{}0 + pi\textasciicircum{}0
\end{definition}
\begin{proof}
  This is the declared model definition of particle Side.
\end{proof}
\begin{definition}[Figure Panel]
  \label{def:physics:phyx-mini-0554:phyxminiproblems-problemphyxmini0554-figurepanel}
  \lean{PhyXMiniProblems.ProblemPhyXMini0554.FigurePanel}
  The before- and after-interaction panels labelled (a) and (b)
\end{definition}
\begin{proof}
  This is the declared model definition of Figure Panel.
\end{proof}
\begin{definition}[Planar Direction]
  \label{def:physics:phyx-mini-0554:phyxminiproblems-problemphyxmini0554-planardirection}
  \lean{PhyXMiniProblems.ProblemPhyXMini0554.PlanarDirection}
  Qualitative momentum-arrow directions visible in the primary image
\end{definition}
\begin{proof}
  This is the declared model definition of Planar Direction.
\end{proof}
\begin{definition}[Kaon Proton Reaction Figure]
  \label{def:physics:phyx-mini-0554:phyxminiproblems-problemphyxmini0554-kaonprotonreactionfigure}
  \lean{PhyXMiniProblems.ProblemPhyXMini0554.KaonProtonReactionFigure}
  \uses{def:physics:phyx-mini-0554:phyxminiproblems-problemphyxmini0554-diagramaxis, def:physics:phyx-mini-0554:phyxminiproblems-problemphyxmini0554-particlelabel, def:physics:phyx-mini-0554:phyxminiproblems-problemphyxmini0554-figurepanel, def:physics:phyx-mini-0554:phyxminiproblems-problemphyxmini0554-planardirection}
  Qualitative labels and geometry transcribed from the primary raster. The two angles are uncalibrated real radian readouts because the image names but does not numerically measure them
\end{definition}
\begin{proof}
  This is the declared model definition of Kaon Proton Reaction Figure.
\end{proof}
\begin{definition}[Particle State]
  \label{def:physics:phyx-mini-0554:phyxminiproblems-problemphyxmini0554-particlestate}
  \lean{PhyXMiniProblems.ProblemPhyXMini0554.ParticleState}
  \uses{def:physics:phyx-mini-0554:phyxminiproblems-problemphyxmini0554-restmassquantity, def:physics:phyx-mini-0554:phyxminiproblems-problemphyxmini0554-planemomentum}
  The independent physical state of one massive particle. In particular, the lambda kinetic energy is an unconstrained dimensionful field here rather than being defined from the recorded answer
\end{definition}
\begin{proof}
  This is the declared model definition of Particle State.
\end{proof}
\begin{definition}[Kaon Proton Reaction Setup]
  \label{def:physics:phyx-mini-0554:phyxminiproblems-problemphyxmini0554-kaonprotonreactionsetup}
  \lean{PhyXMiniProblems.ProblemPhyXMini0554.KaonProtonReactionSetup}
  \uses{def:physics:phyx-mini-0554:phyxminiproblems-problemphyxmini0554-particlelabel, def:physics:phyx-mini-0554:phyxminiproblems-problemphyxmini0554-kaonprotonreactionfigure, def:physics:phyx-mini-0554:phyxminiproblems-problemphyxmini0554-particlestate}
  All four particle states and the supplied two-panel figure
\end{definition}
\begin{proof}
  This is the declared model definition of Kaon Proton Reaction Setup.
\end{proof}
\begin{definition}[total Relativistic Energy In Joules]
  \label{def:physics:phyx-mini-0554:phyxminiproblems-problemphyxmini0554-totalrelativisticenergyinjoules}
  \lean{PhyXMiniProblems.ProblemPhyXMini0554.totalRelativisticEnergyInJoules}
  \uses{def:physics:phyx-mini-0554:phyxminiproblems-problemphyxmini0554-energyinjoules, def:physics:phyx-mini-0554:phyxminiproblems-problemphyxmini0554-restenergyequivalentinjoules, def:physics:phyx-mini-0554:phyxminiproblems-problemphyxmini0554-particlestate}
  Total relativistic energy m c\textasciicircum{}2 + K of a particle, in joules
\end{definition}
\begin{proof}
  This is the declared model definition of total Relativistic Energy In Joules.
\end{proof}
\begin{definition}[Matches Reaction Problem Data]
  \label{def:physics:phyx-mini-0554:phyxminiproblems-problemphyxmini0554-matchesreactionproblemdata}
  \lean{PhyXMiniProblems.ProblemPhyXMini0554.MatchesReactionProblemData}
  \uses{def:physics:phyx-mini-0554:phyxminiproblems-problemphyxmini0554-diagramaxis, def:physics:phyx-mini-0554:phyxminiproblems-problemphyxmini0554-energyinjoules, def:physics:phyx-mini-0554:phyxminiproblems-problemphyxmini0554-energyinmegaelectronvolts, def:physics:phyx-mini-0554:phyxminiproblems-problemphyxmini0554-restmassinmegaelectronvoltsperc2, def:physics:phyx-mini-0554:phyxminiproblems-problemphyxmini0554-momentumcomponentinsi, def:physics:phyx-mini-0554:phyxminiproblems-problemphyxmini0554-kaonprotonreactionsetup}
  Numerical quantities stated in the prose. The lambda kinetic energy is deliberately absent. The proton-at-rest condition fixes both its kinetic energy and its two spatial-momentum components to zero
\end{definition}
\begin{proof}
  This is the declared model definition of Matches Reaction Problem Data.
\end{proof}
\begin{definition}[Matches Primary Reaction Figure]
  \label{def:physics:phyx-mini-0554:phyxminiproblems-problemphyxmini0554-matchesprimaryreactionfigure}
  \lean{PhyXMiniProblems.ProblemPhyXMini0554.MatchesPrimaryReactionFigure}
  \uses{def:physics:phyx-mini-0554:phyxminiproblems-problemphyxmini0554-momentumcomponentinsi, def:physics:phyx-mini-0554:phyxminiproblems-problemphyxmini0554-momentummagnitudeinsi, def:physics:phyx-mini-0554:phyxminiproblems-problemphyxmini0554-kaonprotonreactionsetup}
  Primary-image evidence. The outgoing component formulas express the two arrow directions at the named acute angles. No angle value and no lambda kinetic-energy value is inferred from the raster
\end{definition}
\begin{proof}
  This is the declared model definition of Matches Primary Reaction Figure.
\end{proof}
\begin{definition}[Has Physical Reaction Parameters]
  \label{def:physics:phyx-mini-0554:phyxminiproblems-problemphyxmini0554-hasphysicalreactionparameters}
  \lean{PhyXMiniProblems.ProblemPhyXMini0554.HasPhysicalReactionParameters}
  \uses{def:physics:phyx-mini-0554:phyxminiproblems-problemphyxmini0554-restmassinkilograms, def:physics:phyx-mini-0554:phyxminiproblems-problemphyxmini0554-energyinjoules, def:physics:phyx-mini-0554:phyxminiproblems-problemphyxmini0554-momentummagnitudeinsi, def:physics:phyx-mini-0554:phyxminiproblems-problemphyxmini0554-kaonprotonreactionsetup}
  Positivity conditions selecting physically meaningful particle states
\end{definition}
\begin{proof}
  This is the declared model definition of Has Physical Reaction Parameters.
\end{proof}
\begin{definition}[Satisfies Relativistic Reaction Laws]
  \label{def:physics:phyx-mini-0554:phyxminiproblems-problemphyxmini0554-satisfiesrelativisticreactionlaws}
  \lean{PhyXMiniProblems.ProblemPhyXMini0554.SatisfiesRelativisticReactionLaws}
  \uses{def:physics:phyx-mini-0554:phyxminiproblems-problemphyxmini0554-diagramaxis, def:physics:phyx-mini-0554:phyxminiproblems-problemphyxmini0554-speedoflightinmeterspersecond, def:physics:phyx-mini-0554:phyxminiproblems-problemphyxmini0554-restenergyequivalentinjoules, def:physics:phyx-mini-0554:phyxminiproblems-problemphyxmini0554-momentumcomponentinsi, def:physics:phyx-mini-0554:phyxminiproblems-problemphyxmini0554-momentummagnitudeinsi, def:physics:phyx-mini-0554:phyxminiproblems-problemphyxmini0554-particlelabel, def:physics:phyx-mini-0554:phyxminiproblems-problemphyxmini0554-kaonprotonreactionsetup, def:physics:phyx-mini-0554:phyxminiproblems-problemphyxmini0554-totalrelativisticenergyinjoules}
  Governing laws for the isolated two-body reaction: each massive particle obeys E\textasciicircum{}2 = (m c\textasciicircum{}2)\textasciicircum{}2 + (p c)\textasciicircum{}2, with E = m c\textasciicircum{}2 + K; total relativistic energy is conserved; and both displayed spatial-momentum components are conserved. These laws relate independent state fields and contain no numerical lambda kinetic energy or displayed answer choice
\end{definition}
\begin{proof}
  This is the declared model definition of Satisfies Relativistic Reaction Laws.
\end{proof}
\begin{definition}[Answer Choice]
  \label{def:physics:phyx-mini-0554:phyxminiproblems-problemphyxmini0554-answerchoice}
  \lean{PhyXMiniProblems.ProblemPhyXMini0554.AnswerChoice}
  Labels attached to the four answer choices
\end{definition}
\begin{proof}
  This is the declared model definition of Answer Choice.
\end{proof}
\begin{definition}[displayed Lambda Kinetic Energy Me V]
  \label{def:physics:phyx-mini-0554:phyxminiproblems-problemphyxmini0554-displayedlambdakineticenergymev}
  \lean{PhyXMiniProblems.ProblemPhyXMini0554.displayedLambdaKineticEnergyMeV}
  \uses{def:physics:phyx-mini-0554:phyxminiproblems-problemphyxmini0554-answerchoice}
  Kinetic-energy readout in MeV printed beside each answer choice
\end{definition}
\begin{proof}
  This is the declared model definition of displayed Lambda Kinetic Energy Me V.
\end{proof}
\begin{definition}[recorded Answer Choice]
  \label{def:physics:phyx-mini-0554:phyxminiproblems-problemphyxmini0554-recordedanswerchoice}
  \lean{PhyXMiniProblems.ProblemPhyXMini0554.recordedAnswerChoice}
  \uses{def:physics:phyx-mini-0554:phyxminiproblems-problemphyxmini0554-answerchoice}
  Dataset metadata records answer choice D
\end{definition}
\begin{proof}
  This is the declared model definition of recorded Answer Choice.
\end{proof}
% --- Archon declaration topology end ---
% --- Archon physics formalization source end ---
